\begin{frame}{Ratios de liquidez}
\color{liquidez}\textbf{Capacidad para atender obligaciones de corto plazo}\color{black}

\vspace{0.7em}
Estos ratios relacionan activos corrientes con pasivos corrientes. Su análisis debe revisar la cobrabilidad de las cuentas por cobrar y la velocidad de realización de inventarios.
\end{frame}

\begin{frame}{Razón corriente}
\ratio{Qué evalúa}{Cobertura general de deudas corrientes con activos corrientes.}
\vspace{0.8em}
\[
\text{Razón corriente}=\frac{\text{Activo corriente}}{\text{Pasivo corriente}}
\]
\textbf{Datos necesarios:} activo corriente y pasivo corriente.
\end{frame}

\begin{frame}{Caso práctico: razón corriente}
\casoeducativo

\vspace{0.4em}
Comercial Andina S.A.C. presenta activo corriente de \soles{480,000.00} y pasivo corriente de \soles{300,000.00}.
\[
\frac{480,000.00}{300,000.00}=1.60
\]
\textbf{Interpretación:} por cada \soles{1.00} de obligación corriente, dispone de \soles{1.60} en activos corrientes.

\textbf{Decisión:} revisar composición del activo corriente antes de asumir que la liquidez es suficiente.
\end{frame}

\begin{frame}{Prueba ácida}
\ratio{Qué evalúa}{Cobertura inmediata sin depender de la venta de inventarios.}
\[
\text{Prueba ácida}=\frac{\text{Activo corriente}-\text{Inventarios}}{\text{Pasivo corriente}}
\]
\textbf{Datos necesarios:} activo corriente, inventarios y pasivo corriente.
\end{frame}

\begin{frame}{Caso práctico: prueba ácida}
\casoeducativo
\[
\frac{480,000.00-180,000.00}{300,000.00}=1.00
\]
\textbf{Interpretación:} sin inventarios, la empresa posee \soles{1.00} de activos líquidos por cada \soles{1.00} de deuda corriente.

\textbf{Decisión:} sostener caja y cobranzas, porque la cobertura ajustada no deja amplio margen.
\end{frame}

\begin{frame}{Comparación de liquidez}
\begin{center}
\begin{tabular}{lcc}
\toprule
Indicador & Resultado & Lectura \\
\midrule
Razón corriente & 1.60 & Cobertura general \\
Prueba ácida & 1.00 & Cobertura sin inventarios \\
\bottomrule
\end{tabular}
\end{center}
\vspace{0.5em}
La diferencia de 0.60 muestra dependencia de inventarios. La decisión no se basa solo en el valor alto, sino en la rapidez con que esos inventarios pueden convertirse en efectivo.
\end{frame}
